\documentclass{article}
\usepackage[french]{babel}
\usepackage[utf8]{inputenc}
\usepackage[T1]{fontenc}
\usepackage{amsmath}
\usepackage{amsfonts}
\usepackage{enumitem}
\usepackage{xcolor}
\usepackage{booktabs}

\begin{document}

\section{MultiAgentMADDPG\_\-LatencyEnergySum \\ Pilier central du projet}

MultiAgentMADDPG\_LatencyEnergySum constitue l’un des cœurs algorithmiques du projet.  
Son objectif est de résoudre le problème du \textbf{Split Inference} (partitionnement dynamique de modèles d’intelligence artificielle) dans un réseau IoT hétérogène, en s’appuyant sur l’\textbf{apprentissage par renforcement multi-agent}.

\subsection{Objectif principal --- Fonction de récompense multi-objectif}

Contrairement aux versions antérieures focalisées uniquement sur la latence ou sur des contraintes binaires (succès/échec), cette implémentation adopte une optimisation multi-objectif explicite :

\begin{equation}
    r = -(\text{Latence totale} + \text{Énergie totale})
\end{equation}

\begin{itemize}[label=$\bullet$]
    \item \textbf{Latence totale} = temps de calcul + temps de transmission sur l’ensemble du pipeline d’inférence
    \item \textbf{Énergie totale} = consommation cumulée des appareils, calculée via le modèle physique
    \[
    E = P \times T
    \]
    (Puissance instantanée $\times$ durée d’utilisation)
    \item \textbf{Objectif} : contraindre les agents à trouver des stratégies qui soient \textbf{à la fois rapides} et \textbf{très économes en énergie}, afin de maximiser la durée de vie des batteries des objets IoT.
\end{itemize}

\subsection{Algorithme sous-jacent : MADDPG (CTDE)}

L’algorithme employé est le \textbf{Multi-Agent Deep Deterministic Policy Gradient} (MADDPG), qui repose sur le paradigme CTDE :

\begin{description}[leftmargin=*,font=\normalfont\bfseries]
    \item[Centralized Training] 
    Pendant la phase d’apprentissage, le(s) critique(s) ont une vision globale :
    \begin{itemize}[label=--]
        \item état complet du réseau
        \item actions de \textbf{tous} les agents
    \end{itemize}
    Cela permet de stabiliser l’apprentissage dans un environnement non stationnaire multi-agent.

    \item[Decentralized Execution] 
    Une fois l’entraînement terminé, chaque agent prend ses décisions de façon totalement décentralisée :
    \begin{itemize}[label=--]
        \item ne voit que son observation locale
        \item décide seul du placement des couches qui lui sont confiées
    \end{itemize}
\end{description}

\subsection{Principaux enseignements tirés des résultats}

\begin{enumerate}[label=\arabic*.]
    \item \textbf{Normalisation efficace latence–énergie} \\
    En plaçant latence et énergie sur la même échelle de grandeur, le compromis entre les deux objectifs devient réellement réalisable ; l’agent ne sacrifie plus systématiquement l’un pour l’autre.

    \item \textbf{Adaptabilité automatique} \\
    Le système détecte et exploite spontanément les appareils plus performants, plus économes ou plus sécurisés sans règle explicite.
\end{enumerate}

\subsection{Synthèse}

MultiAgentMADDPG\_LatencyEnergySum représente la version la plus aboutie et réaliste du système développé.  
Il parvient à gérer efficacement le \textbf{trade-off fondamental} entre :

\begin{center}
    \large
    \textbf{vitesse d’inférence} \quad $\longleftrightarrow$ \quad \textbf{survie énergétique du réseau IoT contraint}
\end{center}

\bigskip

C’est actuellement le pilier algorithmique le plus robuste et le plus prometteur pour un déploiement futur sur de vrais dispositifs hétérogènes.

\subsection{1. Agents et modèles d'inférence}

Trois agents sont entraînés simultanément, chacun étant responsable du partitionnement d’un modèle d’intelligence artificielle de complexité différente.

\begin{table}[ht]
\centering
\renewcommand{\arraystretch}{1.3}
\begin{tabular}{>{\ttfamily}l l c l}
\toprule
Agent & Modèle & Nb couches & Description \\
\midrule
0     & hugcnn     & 10  & Très lourd ($\sim$520 Mo RAM pour certaines couches) \\
1     & cnn15      & 15  & Couches très légères ($\sim$4 Mo chacune) \\
2     & miniresnet & 5   & Taille moyenne, mais gourmand en calcul \\
\bottomrule
\end{tabular}
\caption{Description des trois agents et de leurs modèles associés}
\label{tab:agents}
\end{table}

\subsection{2. Infrastructure du réseau IoT (5 appareils hétérogènes)}

\begin{table}[ht]
\centering
\small
\renewcommand{\arraystretch}{1.25}
\begin{tabular}{clllll}
\toprule
ID & Vitesse CPU & RAM & Bande passante & Score de confiance \\
\midrule
0  & $\sim$51.4  & 817      & 205                   & 0.83             \\
1  & $\sim$46.4  & 871      & 298                   & 0.94 \\
2  & $\sim$50.9  & 822      & 168                   & 0.54   \\
3  & $\sim$50.0  & 818      & 190                   & 0.71    \\
4  & $\sim$47.2  & 1212     & 312                   & 0.51  \\
\bottomrule
\end{tabular}
\caption{Caractéristiques des 5 appareils IoT hétérogènes}
\label{tab:iot-devices}
\end{table}

\subsection{3. Principaux paramètres de l'environnement}

\begin{itemize}[label=$\bullet$]
    \item \textbf{Budget énergétique initial} : batterie aléatoire entre 500 et 1200 unités par appareil
    \item \textbf{Modèle de consommation énergétique} :
    \[
    E = P \times T \qquad \text{avec} \quad
    \begin{cases}
        \alpha_{\text{comp}} = 1.0~\text{W (calcul)} \\
        \alpha_{\text{comm}} = 1.0~\text{W (communication)}
    \end{cases}
    \]
    \item \textbf{Niveaux de confidentialité} : 0 (public) $\rightarrow$ 3 (top secret)
    \begin{itemize}[label=--]
        \item Niveau 3 $\implies$ Trust $\ge 0.80$ $\quad\Rightarrow$ seuls appareils 0 et 1 sont éligibles
    \end{itemize}
\end{itemize}

\subsection{Synthèse -- Principaux conflits intentionnellement créés}

\begin{enumerate}[label=\arabic*.]
    \item \textbf{Conflit de sécurité / confiance} \\
    Les trois agents ont initialement tendance à converger vers les appareils 0 et 1 (seuls éligibles pour les couches de niveau 3).

    \item \textbf{Conflit mémoire (split forcé)} \\
    Le modèle \texttt{hugcnn} (agent 0) dépasse largement la capacité RAM d’un seul appareil $\rightarrow$ oblige à trouver des points de coupure (split points) intelligents.

    \item \textbf{Conflit énergétique et partage de ressources} \\
    L’appareil le plus capacitif en RAM (n°4) n’est pas celui qui a la plus grande batterie $\rightarrow$ pénalise la stratégie « tout sur le plus gros ».
\end{enumerate}

\bigskip

Cette architecture reproduit de manière réaliste les principaux trade-offs rencontrés dans un déploiement de \textbf{Split Inference} sur un réseau IoT contraint et hétérogène.

\end{document}